\section{Introduction}\label{sec:intro}

Discharge summaries of intensive care unit (ICU) stays condense the clinical history,
organ support, infections, complications and outcome of a critically ill patient into
narrative text. Converting that narrative into the variables of a clinical registry or a
research database still depends largely on manual chart abstraction, which is slow, costly
and difficult to scale beyond small cohorts. Large language models (LLMs) have shown that
clinical information can be extracted from free text with little or no task-specific
training~\cite{agrawal2022fewshot}, and a growing body of work applies them to structured
extraction from clinical notes, pathology reports and radiology reports, including with open-weight
models run on local infrastructure%
~\cite{privacy2024structured,iterativerefine2025,clinicalentityretrieval2024,spaanderman2025structured}.

For research use, the value of an extracted variable depends on whether it can be audited.
A reviewer must be able to see which part of the note supports each answer, to distinguish
a documented absence from a missing answer, and to detect when the output violates the
structure of the form. These requirements become harder to meet as the form grows. When a
single generation must return hundreds of variables, the output becomes long: it may exceed
the token budget and be truncated, it may omit keys or break the JSON structure, and an
error anywhere invalidates or obscures the whole response. A common engineering response is
to decompose the task, for example by processing the note sentence by sentence, by reviewing
and correcting an initial output, by chaining sequential sub-prompts, or by retrieving only the
passages relevant to each variable~\cite{llmie2025,fornasiere2024medical,spaanderman2025structured,clinicalentityretrieval2024}. Decomposition is not
free, however. It multiplies the number of calls, repeats the note in every prompt, increases
latency and creates new opportunities for inconsistency between answers produced
independently for related variables.

Existing studies differ in note type, language, schema size, annotation procedure, models and
metrics, and they rarely isolate the effect of the calling strategy on the same patients, the
same schema and the same models (Section~\ref{sec:related}). In information extraction from
radiology reports, a scoping review found that results could not be compared across studies and
that annotation guidelines, inter-annotator agreement and code were often not
reported~\cite{scoping2024radiology}. The choice between a
single call and a decomposed pipeline is therefore usually made on engineering intuition rather
than on a paired measurement of what each strategy gains and costs.

This study addresses that gap for one concrete and demanding case: a hierarchical annotation form
of 199 variables applied to Spanish discharge summaries of an adult ICU in a public hospital in Chile, the
same form used by human annotators in an annotation platform. We compare three strategies that share the
note, the rules, the output contract and the validator, and differ only in how the form is requested: one
call for all 199 variables; a schema-constrained harness of 18 independent calls, each returning a clinically
coherent group of 2 to 16 variables that is validated, retried on error and merged; and the same harness in
which a preliminary call first builds an evidence-grounded model of the case (ICU stays, timeline, problem
status, support, infections and outcome) that every group call receives. The third strategy addresses a
limitation of independent calls: many variables require inference across the whole episode, and groups that
cannot see each other's answers may reconstruct the clinical story differently. In all strategies the model
cites evidence by line identifier and the software copies the cited text, so that a quotation cannot be
paraphrased or invented by the model.

The contribution is not the use of several prompts, nor the idea of an intermediate representation, both of
which have precedents. It is (i) a controlled, within-patient and within-model comparison of monolithic,
schema-decomposed and case-model-guided extraction under an identical 199-variable contract; (ii) a
reproducible protocol that separates structural validity, evidence fidelity, clinical coherence, agreement
with human annotators and operational cost, and that never converts a failed or truncated answer into a
negative finding; and (iii) an execution with three open-weight models on local infrastructure, with frozen
inputs, hashes and per-call traces.
